\chapter{Машина в цілому}
% full version: chapters/10_synthesis.tex

\pencilfig{10}{\pencilart{10}}{Машина в цілому: шина даних, гальма в ній і чотири радіоканали керування згори.}

Ми брали по одному типу нейронів і питали, що він додає до машини. Час зібрати деталі разом.

Вийшла машина з трьох шарів. Унизу --- величезна телефонна мережа \nt{GLUT}-нейронів: нею йдуть дані. Поруч --- гальмівні \nt{GABA}-нейрони: вони керують потоком, вирішують, чому пройти, коли і наскільки гучно. А над усім цим --- кілька радіостанцій, у кожної своя ручка. \nt{DOPA} каже, які зміни зв'язків залишити. \nt{SERO} тримає загальний рівень і, за гіпотезою, горизонт терпіння. \nt{NORA} обирає між пошуком і рішенням. \nt{ACET} --- між записом і читанням. Кожна радіостанція насправді --- невеликий пучок каналів, а не один провід, але каналів однаково жменька на вісімдесят шість мільярдів нейронів.

\section*{Одна історія}

Простежмо одну подію через усю машину. Тварина давно знає: натиснеш важіль A --- отримаєш сік. Одного дня світ змінюється: A більше не працює, а сік приносить важіль B, яким тварина майже не користувалася.

Сік після A не прийшов --- дофамін відповідає провалом, і зв'язки, що обрали A, слабшають. Провали повторюються, помилки стають більшими, ніж зазвичай, --- норадреналін, за гіпотезою, повідомляє: світ змінився. Ручка контрастності опускається, вибір стає м'якшим, і шум частіше підштовхує спробувати B. Ацетилхолін переводить мережу в режим запису. Спроба B приносить несподіваний сік --- дофамін спалахує, зв'язки, що обрали B, міцнішають. За кілька спроб очікування наздоганяє новий світ, сюрпризи закінчуються, ручки повертаються на місце.

Зауважте: жодна ручка не знає, що роблять інші. Кожна відгукується на свої ознаки, і порядок дій складається сам, як в асинхронній схемі, де кожен блок спрацьовує, коли готові його входи. Що кожна ручка поводиться так окремо, здебільшого виміряно. Що вони злагоджено працюють разом --- поки гіпотеза.

\section*{Чим вона не схожа на ваш ноутбук}

Ноутбук має спільний такт --- мозок ні. Елементи ноутбука надійні --- синапс губить понад половину сигналів. Пам'ять ноутбука окремо від процесора --- у мозку вона в тих самих зв'язках. Ноутбук учиться (якщо вчиться) зворотним поширенням помилки --- мозок місцевими правилами й одним широкомовним «краще, ніж очікувалося». І головне: ноутбук витрачає десятки ватів на один процесор, а мозок --- двадцять ватів на все.

\section*{Чого ми не знаємо}

Ми не знаємо, наскільки кора користується точним часом клацань. Не знаємо, як вона налаштовує мільярди зв'язків у багатоповерхових колах, якщо вчитель у неї один на всіх. Не до кінця розуміємо, що передають радіостанції. Не знаємо, як далекі частини мозку узгоджують роботу без спільного такту. І поки ніхто не вміє побудувати машину, яка робила б те саме на двадцяти ватах.

Далі книга виходить за межі цього ядра: до входів і виходів машини, до її налагодження, до сну й до машин із живих нейронів на чипі.

\fullversion{10}
